\documentclass[header={Ancillary Note: Bounded vs Continuous},notheorems]{study-note}

\begin{document}

\StudyTitleBlock{N5}{Ancillary Note}{Bounded vs. continuous linear maps}

\vspace{0.8em}

\begin{focusbox}
One of the first recurring facts in functional analysis is that for linear maps between normed spaces, boundedness and continuity are equivalent. This note isolates that fact clearly.
\end{focusbox}

\tableofcontents

\section{Definitions}

\begin{defbox}{Bounded linear map}
Let $T:X\to Y$ be linear between normed spaces. Then $T$ is \textbf{bounded} if there exists $C\ge 0$ such that
\[
\norm{Tx}_Y\le C\norm{x}_X
\qquad \forall x\in X.
\]
\end{defbox}

\begin{defbox}{Continuous map at a point}
A map $T:X\to Y$ is \textbf{continuous at} $x_0\in X$ if for every $\varepsilon>0$ there exists $\delta>0$ such that
\[
\norm{x-x_0}_X<\delta
\Longrightarrow
\norm{Tx-Tx_0}_Y<\varepsilon.
\]
\end{defbox}

\section{Main Equivalence}

\begin{thmbox}{Boundedness equals continuity for linear maps}
Let $T:X\to Y$ be linear between normed spaces. Then the following are equivalent:
\begin{enumerate}
\item $T$ is bounded;
\item $T$ is continuous everywhere;
\item $T$ is continuous at one point;
\item $T$ is continuous at $0$.
\end{enumerate}
\end{thmbox}

\paragraph{Why continuity at one point is enough.}
Linearity lets you shift everything back to the origin. Once you understand continuity at $0$, you understand continuity everywhere.

\begin{exbox}{Proof sketch: bounded implies continuous}
If
\[
\norm{Tx}\le C\norm{x},
\]
then
\[
\norm{Tx-Ty}=\norm{T(x-y)}\le C\norm{x-y}.
\]
So $T$ is even Lipschitz continuous.
\end{exbox}

\begin{exbox}{Proof sketch: continuity at $0$ implies bounded}
If $T$ is continuous at $0$, then there exists $\delta>0$ such that
\[
\norm{x}<\delta \Longrightarrow \norm{Tx}<1.
\]
For $x\neq 0$, apply this to
\[
z=\frac{\delta}{2\norm{x}}x.
\]
Then $\norm{z}<\delta$, so $\norm{Tz}<1$, which gives
\[
\norm{Tx}<\frac{2}{\delta}\norm{x}.
\]
Thus $T$ is bounded.
\end{exbox}

\section{Operator Norm}

\begin{defbox}{Operator norm}
For a bounded linear map $T:X\to Y$, define
\[
\norm{T}_{\mathcal L(X,Y)}:=\sup_{\norm{x}_X\le 1}\norm{Tx}_Y.
\]
\end{defbox}

\begin{thmbox}{Equivalent formula}
For $T\neq 0$,
\[
\norm{T}=\sup_{x\neq 0}\frac{\norm{Tx}}{\norm{x}}.
\]
\end{thmbox}

\section{What To Remember Quickly}

\begin{focusbox}
Quick summary:
\begin{enumerate}
\item for linear maps, bounded = continuous;
\item it is enough to check continuity at $0$;
\item bounded linear maps are automatically Lipschitz;
\item the operator norm is the best constant in the inequality
\[
\norm{Tx}\le C\norm{x}.
\]
\end{enumerate}
\end{focusbox}

\end{document}
